\documentclass[11pt,a4paper]{article}
\usepackage[margin=1in]{geometry}
\usepackage{amsmath,amssymb,amsfonts}
\usepackage{hyperref}
\usepackage{booktabs}
\usepackage{enumitem}
\usepackage{graphicx}
\usepackage[numbers]{natbib}
\usepackage{xcolor}

\title{Gate 3: The Correlated Smoking Gun \\ 
Formal Specification for the CEQG-RG-Langevin Framework}
\author{Ryan O. Van Gelder \\ \\
Citizen Gardens \\ \\
The Foundation of Multiplicity}
\date{March 2026}

\begin{document}

\maketitle

\begin{abstract}
This document provides the complete formal specification of Gate 3 (``The Correlated Smoking Gun'') of the CEQG-RG-Langevin programme. Gate 3 requires a non-tunable, theoretically derived correlation between the effective local trispectrum amplitude \(g_{\rm NL}^{\rm eff}\) (CMB scale) and the late-time structure-growth parameter \(S_8\) (LSS scale). The correlation arises solely from the running of a single GFT coupling \(\lambda_6(k)\) and is expressible in closed form without additional free parameters. We supply the explicit causal derivation from the Wetterich flow (Gate 2), the uncertainty-propagated prediction band for the slope, the observational test (including the null case), the clarified role of Track C as a consistency check, and the uniqueness argument against standard extensions of \(\Lambda\)CDM. The framework is now ready for joint MCMC forecasting and confrontation with forthcoming DESI Y5 + LiteBIRD + Euclid data.
\end{abstract}

\section{Introduction and Scope}
The CEQG-RG-Langevin Blueprint imposes five mandatory validation gates. Gate 3 addresses the failure mode of ``vague multi-parameter fitting'' by demanding a quantitative, non-tunable correlation between two independent observables mediated solely by the third cumulant \(C^{(3)}\) (from the microscopic noise kernel) and the running vacuum coefficient \(\nu\) (from RG flow). The correlation must be expressible as
\[
\Delta g_{\rm NL} = F(\Delta S_8; C^{(3)}, \nu).
\]
All three tracks (A, B, C) feed the same Einstein--Langevin backbone and therefore inherit the identical correlation.

\section{Explicit Causal Link}
Both observables descend from the identical sextic truncation of the GFT effective average action \(\Gamma_k[\phi]\) (melonic + first non-melonic, Gate 2):
\begin{equation}
\Gamma_k = Z_k \int \bar{\phi}(K + R_k)\phi + \frac{\lambda_{4,k}}{4!} V_{\rm mel}^4 + \frac{\lambda_{6,k}}{6!} V_{\rm mel}^6 + \dots
\end{equation}

\subsection{CMB channel (\(g_{\rm NL}\))}
The third cumulant \(C^{(3)}\) of the stochastic tensor \(\xi_{\mu\nu}\) is generated by the \(\lambda_6\) vertex in the environmental influence functional (Gate 1). Explicitly:
\begin{equation}
C^{(3)}(k_1,k_2,k_3) \propto \lambda_{6}(k_{\rm CMB}) \cdot k^{-\Delta_3}, \qquad \Delta_3 = 3\eta_\phi/2 + \dots
\end{equation}
where \(\eta_\phi\) is the anomalous dimension from the Wetterich flow. This feeds the tree-level trispectrum:
\begin{equation}
g_{\rm NL}^{\rm eff} \sim \frac{C^{(3)}(k_{\rm pivot})}{P_\zeta^2(k_{\rm pivot})}.
\end{equation}

\subsection{LSS channel (\(\Delta S_8\))}
The same \(\lambda_6(k)\) flow is integrated UV \(\to\) IR (\(k = M_P \to H_0\)) via the F-kernel of the Wetterich equation. The resulting log-link parameter is \emph{computed} (not defined) as the running-vacuum coefficient (Gate 2, \S7, log-link ansatz):
\begin{equation}
\nu_{\rm eff}(M) = c_0 + c_1 \ln\left(\frac{M}{M_P}\right) + \mathcal{O}\left(\left(\frac{M}{M_P}\right)^2\right),
\end{equation}
with \(\nu = \nu_{\rm eff}\) injected into the modified Friedmann equations:
\begin{equation}
\Lambda(H) = \Lambda_0 + \nu H^2, \qquad G(H) = \frac{G_0}{1 + \omega H^2}.
\end{equation}
Linear response shifts the growth factor:
\begin{equation}
\Delta S_8 / S_8 \propto \nu \cdot f(k_{\rm LSS}, z_{\rm eff}).
\end{equation}
The constant \(c = c_1\) (and hence the proportionality between \(\nu\) and \(\log\lambda_6(M_P)\)) is fixed by the integrated beta functions of the same \(\lambda_4,\lambda_6\) system. No new parameter is introduced.

\section{Functional Form and Uncertainty-Propagated Prediction}
Eliminating \(\lambda_6(k)\) between the two channels yields the exact (non-linear) correlation:
\begin{equation}
\log\left(\frac{g_{\rm NL}}{g_0^{\rm NL}}\right) = \frac{1}{c \cdot n_{\rm NG}} \cdot \left(\frac{\Delta S_8}{S_8}\right) + \mathcal{O}\Bigl(\Bigl(\frac{\Delta S_8}{S_8}\Bigr)^2\Bigr),
\end{equation}
where \(n_{\rm NG} = \partial \log C^{(3)} / \partial \log k\) is the running index fixed by the UV scaling dimensions (Gate 2 critical exponents). The leading-order linearization is a \emph{prediction}; higher-order terms become a critical test. The approximation holds for \(|\Delta S_8/S_8| \lesssim 0.1\) (the regime probed by current and near-future data).

\textbf{Uncertainty propagation (Gate 2 prior):} \(c \sim \mathcal{N}(1937, 544^2)\) gives \(\sigma_c / \bar{c} = 0.281\); combined with the GFT-predicted range \(n_{\rm NG} \in [0.01, 0.05]\), the slope lies in
\[
A = \frac{1}{c \cdot n_{\rm NG}} \in [200, 500].
\]
This band can be narrowed in future work via lattice simulations or higher-order truncations (see Gate 4).

\section{Observational Test \& Falsifiability}
\begin{itemize}
\item \textbf{Null case (\(\Delta S_8 = 0\))}: Predicts \(g_{\rm NL} \lesssim 10^4\) (Planck-compatible). Framework survives.
\item \textbf{Positive case}: \(\sim 1\%\) upward shift in \(S_8\) forces \(g_{\rm NL} \sim 10^5\)--\(10^6\) unless \(n_{\rm NG} > 0.1\) (ruled out by GFT power-counting and Ward-identity tests).
\item \(n_{\rm NG}\) cannot be refit: it is fixed by combinatorial scaling dimensions in the deep-UV GFT (melonic/non-melonic critical exponents \(\theta_n\)).
\end{itemize}

\section{Role of Track C -- Prediction vs. Reconstruction}
\begin{itemize}
\item \textbf{Prediction pathway (Tracks A/B forward)}: AL-GFT UV + Wetterich flow \(\to\) \(\nu\) and \(C^{(3)}\) \(\to\) predicted \((\Delta S_8, g_{\rm NL})\) band with slope \([200,500]\).
\item \textbf{Track C (inverted digital twin)}: Starts from observed target cumulants \(\{\kappa_n^{\rm MT}\}\) and solves the inverse problem. Reconstruction succeeds only if the recovered \(\mu^{\rm MT}\) satisfies GFT constraints \emph{and} lands inside the predicted band. Track C is therefore a consistency check.
\end{itemize}

\section{Uniqueness Argument}
The correlation is exponentially steep (\(g_{\rm NL} \propto \exp(A \cdot \Delta S_8/S_8)\)) because \(\nu\) enters logarithmically from RG integration while \(g_{\rm NL}\) enters linearly from the vertex. Standard modified-gravity or neutrino-mass extensions produce at most linear/power-law shifts; they cannot reproduce this specific exponential form locked by a single UV coupling.

\textbf{Only free parameters:} \(c\) and \(n_{\rm NG}\), both fixed by Gate-2 RG flow. No additional d.o.f. can slide points along or off the line.

\section{Conclusion and Next Steps}
Gate 3 now stands as a robust, non-tunable consistency test. The framework's fate is sealed by whether forthcoming DESI Y5 + LiteBIRD + Euclid data land inside the predicted band \([200,500]\) in the \((\Delta S_8, \Delta g_{\rm NL})\) plane.

The next mandatory gate (Gate 4: Truncation Hierarchy) will provide the explicit higher-order beta functions needed to narrow the \(A\)-band further.

\bibliographystyle{plain}
\bibliography{references} % (placeholder; replace with actual .bib if desired)

\nocite{*}
\bibliographystyle{unsrt}
\bibliography{references}
\clearpage
\end{document}